\documentclass[11pt]{article}
\usepackage[margin=1in]{geometry}
\usepackage{amsmath,amssymb,amsthm}
\usepackage[hidelinks]{hyperref}
\newtheorem{theorem}{Theorem}
\title{The two-variable Bernoulli determinant\\
\large A corrected evaluation for MathDB 341858}
\author{Principia Math In Progress}
\date{18 September 2026}
\begin{document}
\maketitle
\begin{abstract}
We prove the intended all-order factorization by integer rank
deficiencies and a leading-coefficient calculation, correcting its
lower product index from zero to one. The underlying evaluation is
also covered by Dominici's published Hahn determinant theorem.
This expository resolution receives zero new-solution credit.
\end{abstract}

\section*{Task 1: the original problem and its resolution}
Normalize Bernoulli polynomials by $B_1(X)=X-\tfrac12$.
For $r\geq1$, the columns of the source's $2r$ by $2r$ matrix
$A_r(x,y)$ are indexed in order by
$(m,v)=(0,1),(0,2),\ldots,(r-1,1),(r-1,2)$.
Put $x_1=x,x_2=y$. Its entries are
\[
 (A_r)_{a,(m,v)}=\delta_{a,m}\quad(0\leq a<r),\qquad
 (A_r)_{r+i,(m,v)}=\frac{B_{i+m+1}(x_v)}{i+m+1}
                 \quad(0\leq i<r).
\]
Then $\bar F_{r,2}(x,y)=\det A_r(x,y)$ is the polynomial
in \cite[(4.7)]{C}. MathDB repeats Conjecture 4.4 of the preprint,
numbered Conjecture 4.1 in the published article.
Its product starts at $j=0$. Literally this is false: at $r=1$
the determinant is $y-x$, whereas the printed expression is $(y-x)^3$.
For example $(x,y)=(0,2)$ gives $2\ne8$.
Here is the complete corrected identity.

\begin{theorem}
For every $r\geq1$ and all complex $x,y$, with $t=y-x$,
\begin{equation}\label{eq:answer}
 \bar F_{r,2}(x,y)=(-1)^{\binom r2}C_r\,t^r
                    \prod_{j=1}^{r-1}(t^2-j^2)^{r-j},
 \qquad
 C_r=\frac{\bigl(\prod_{j=1}^{r-1}j!\bigr)^3}
                 {\prod_{j=r}^{2r-1}j!}.
\end{equation}
Empty products have value one.
\end{theorem}
\begin{proof}
Subtract each column at $x$ from its partner at $y$, then move all
the $x$ columns before all the difference columns. This permutation
has $\binom r2$ inversions. Expanding along the first $r$ rows gives
\begin{equation}\label{eq:reduction}
 \bar F_{r,2}(x,x+t)=(-1)^{\binom r2}t^r
             \det_{0\leq i,j<r}m_{i+j}(x,t),\qquad
 m_k(x,t)=\frac{B_{k+1}(x+t)-B_{k+1}(x)}{(k+1)t}.
\end{equation}
This quotient is polynomial, including at $t=0$.
The identity $B_{k+1}(u+1)-B_{k+1}(u)=(k+1)u^k$ gives
\[
 m_k(x,L)=\frac1L\sum_{a=0}^{L-1}(x+a)^k
 \qquad(L\text{ a positive integer}).
\]
The binomial change of basis from $a^i$ to $(x+a)^i$ is a unit
lower triangular matrix $P(x)$. The moment matrices therefore
satisfy $H(x,L)=P(x)H(0,L)P(x)^T$, so their determinants agree.
Both determinants are polynomials in $x,t$; equality for infinitely
many integer $t=L$ proves equality identically. It remains to evaluate
\[
 \Delta_r(t)=\det_{0\leq i,j<r}\mu_{i+j}(t),\qquad
 \mu_k(t)=\frac{B_{k+1}(t)-B_{k+1}(0)}{(k+1)t}.
\]
The degree of $\mu_k$ is $k$, with leading coefficient $1/(k+1)$.
Every determinant term has degree at most $r(r-1)$, and its
coefficient of that degree is the Hilbert determinant
\begin{equation}\label{eq:Hilbert}
 \det_{0\leq i,j<r}\frac1{i+j+1}
 =\frac{\prod_{0\leq i<j<r}(j-i)^2}
         {\prod_{0\leq i,j<r}(i+j+1)}
 =C_r>0,
\end{equation}
by the Cauchy determinant formula.

For $1\leq s<r$, the matrix at $t=s$ is a moment matrix on $s$
atoms, so has rank at most $s$. This also holds at $t=-s$, since
telescoping in the opposite direction gives
\[
 \mu_k(-s)=\frac1s\sum_{a=-s}^{-1}a^k.
\]
If a polynomial $r$ by $r$ matrix has rank at most $s$ at $t=a$,
an invertible constant row transformation makes $r-s$ rows zero
there. Factoring $t-a$ from each such row proves divisibility of
the determinant by $(t-a)^{r-s}$. Thus
\[
 \prod_{s=1}^{r-1}(t^2-s^2)^{r-s}\mid\Delta_r(t).
\]
The product is monic of degree $r(r-1)$. The degree bound and
\eqref{eq:Hilbert} make the quotient exactly $C_r$.
Equation \eqref{eq:reduction} proves \eqref{eq:answer}, including
all degenerate values by polynomial identity.
\end{proof}

\paragraph{Existing literature and campaign treatment.}
The corrected evaluation receives no new credit. Dominici's
Corollary 4 \cite{D}, combined with his Proposition 1, evaluates
the integer moment determinant. Set both parameters in that
corollary to $1$, put $n=r$, and let $N=L-1$ with $L\geq r+1$.
All weights are one, giving
\begin{align*}
 \det_{0\leq i,j<r}\sum_{a=0}^{L-1}a^{i+j}
 &=\prod_{k=0}^{r-1}
       \frac{(k!)^3(L+k)!}{(L-1-k)!(r+k)!}\\
 &=C_r L^r\prod_{s=1}^{r-1}(L^2-s^2)^{r-s}.
\end{align*}
Translation and polynomial continuation give the full corrected
identity. This identification is our inference from his theorem;
Dominici does not describe it as a solution of Cimpoea\c{s}'s
conjecture. He attributes the summation to earlier work of
Krattenthaler. The elementary rank proof is retained as exposition,
not a claim of a new evaluation. The printed index error earns
no credit either.

\section*{Tasks 2--5: proposed follow-up directions}
These are proposed directions, not claimed solutions; their own
literature checks remain to be done.
\paragraph{Task 2: shifted moment minors.}
Evaluate $\det(\mu_{i+j+q}(t))_{0\leq i,j<r}$ for fixed $q>0$.
Separate forced integer-root factors from the remaining polynomial
and relate the latter to Hahn polynomials.
\paragraph{Task 3: local degeneracy.}
Find explicit polynomial bases for the nullspaces at $t=\pm s$
and the leading nonzero determinant coefficient there.
\paragraph{Task 4: more variables.}
Reconcile the sum/product notation in the source's second
factorization conjecture, then evaluate its intended $r=2$,
arbitrary $D$ determinant, including the rational constant.
\paragraph{Task 5: deformed moments.}
Replace uniform atoms by beta-binomial or $q$-Hahn weights.
Identify surviving rank-based factors and translation properties,
including degenerate parameters.

\paragraph{Verification.}
The exact verifier checks the original block determinant, the
normalized formula, integer ranks, Hilbert constants, and the
published Hahn specialization. These finite checks supplement the
proof. Verification and manuscript review are by the authoring
agent only; no independent peer review is asserted.

\begin{thebibliography}{9}
\bibitem{C}
M. Cimpoea\c{s}, \emph{Determinants with Bernoulli polynomials and the
restricted partition function}, U.P.B. Sci. Bull., Series A
\textbf{81}(4) (2019), 39--52.
\href{https://www.scientificbulletin.upb.ro/static/pdfs/fullf2e_830874.pdf}{Published PDF};
\href{https://arxiv.org/abs/1902.05302}{arXiv:1902.05302}.
\bibitem{D}
D. Dominici, \emph{Hankel determinants, Hahn polynomials, and a formula
of Krattenthaler}, S\'eminaire Lotharingien de Combinatoire
\textbf{81} (2020), Article B81i, 14 pp.
\href{https://www.mat.univie.ac.at/~slc/wpapers/s81dominici.pdf}{Published PDF}.
\end{thebibliography}
\end{document}
